% BHU PYQ -- Manually transcribed & error-corrected
\documentclass[11pt,a4paper]{article}

\usepackage[a4paper,top=2.5cm,bottom=2.5cm,left=2.8cm,right=2.8cm,headheight=14pt]{geometry}
\usepackage{amsmath,amssymb}
\usepackage[T1]{fontenc}
\usepackage[utf8]{inputenc}
\usepackage{lmodern}
\usepackage{microtype}
\usepackage{fancyhdr}
\usepackage{enumitem}
\usepackage{hyperref}
\usepackage{lastpage}
\usepackage{parskip}

\hypersetup{colorlinks=true,urlcolor=black,linkcolor=black,
  pdftitle={CHB-501: Analytical Chemistry-I -- BHU PYQ},pdfauthor={Banaras Hindu University}}

\pagestyle{fancy}\fancyhf{}
\renewcommand{\headrulewidth}{0.4pt}\renewcommand{\footrulewidth}{0.4pt}
\fancyhead[L]{\small   \textit{Banaras Hindu University}}
\fancyhead[R]{\small   \textit{CHB-501: Analytical Chemistry-I}}
\fancyfoot[C]{\small Page  \thepage\ of \pageref{LastPage}}


\newlist{parts}{enumerate}{2}
\setlist[parts,1]{label=(\alph*),leftmargin=*,itemsep=5pt,topsep=3pt}
\setlist[parts,2]{label=(\roman*),leftmargin=*,itemsep=3pt,topsep=2pt}

\newcommand{\pts}[1]{\hfill   \textit{[#1]}}


\begin{document}


\begin{center}
  {\large\bfseries BANARAS HINDU UNIVERSITY}\\[2pt]
  {
\normalsize B.Sc. (Hons.) Semester V Examination, 2016-17}\\[6pt]
  \rule{0.8\linewidth}{0.5pt}\\[6pt]
  {\large\bfseries Paper No. CHB-501: Analytical Chemistry-I}\\[4pt]
  {
\normalsize Time: 3 Hours\hfill Maximum Marks: 70}
\end{center}
\rule{\linewidth}{0.4pt}
\small



\noindent   \textit{Instructions: Answer five questions in total, including Question~1 which is compulsory.}

\normalsize
\bigskip




\noindent   \textbf{Question 1.} Answer all parts of the following: \pts{5 $ \times$ 2 = 10}

\begin{parts}
  \item What is the role of a masking agent? Give one example.
  \item Define systematic error with a suitable example.
  \item Define confidence interval.
  \item Write the chemical structure of oxine (8-hydroxyquinoline).
  \item Differentiate between iodometry and iodimetry.
\end{parts}

\medskip\rule{\linewidth}{0.2pt}




\noindent   \textbf{Question 2.}

\begin{parts}
  \item State the principles of gravimetric analysis. Describe the various steps involved in gravimetric analysis. \pts{8}
  \item Discuss co-precipitation and post-precipitation, highlighting their differences. \pts{6}
\end{parts}

\medskip\rule{\linewidth}{0.2pt}




\noindent   \textbf{Question 3.}

\begin{parts}
  \item Discuss various methods used to minimize errors in quantitative analysis. \pts{6}
  \item The percentage of copper in an alloy was determined as: $12.35\%, 12.42\%, 12.39\%, 12.31\%, 12.44\%$. Calculate: (i) mean, (ii) standard deviation, and (iii) $95\%$ confidence limit. (Given: $t = 2.78$ for 4 degrees of freedom). \pts{8}
\end{parts}

\medskip\rule{\linewidth}{0.2pt}




\noindent   \textbf{Question 4.}

\begin{parts}
  \item Name the titration you would perform for the evaluation of organic functional groups in the following reactions: \pts{8}
  
\begin{parts}
    \item $ \text{R-OH} +  \text{Ac}_2 \text{O}  o  \text{R-OAc} +  \text{AcOH}$
    \item $ \text{R-CHO} +  \text{NH}_2 \text{OH}\cdot \text{HCl}  o  \text{R-CH=NOH} +  \text{H}_2 \text{O} +  \text{HCl}$
  \end{parts}
  \item Discuss the chemical principle, preparation, and standardization of Karl Fischer reagent for the estimation of moisture content. \pts{6}
\end{parts}

\medskip\rule{\linewidth}{0.2pt}




\noindent   \textbf{Question 5.} Discuss the theoretical and practical aspects of the following reagents in chemical analysis: \pts{7 + 7 = 14}

\begin{parts}
  \item EDTA in complexometric titrations
  \item Potassium bromate ($ \text{KBrO}_3$) in redox titrations
\end{parts}

\vfill
\rule{\linewidth}{0.4pt}

\begin{center}\small
  \href{https://bhu-exam.vercel.app/}{Click here to visit our website}\\[4pt]
  \href{https://me-aryan.vercel.app/}{Click here to visit our contributor}\\[4pt]
  \href{https://quashan-me.vercel.app/}{Click here to visit our contributor}
\end{center}

\end{document}
